\section{Introducción}
El presente informe corresponde a la \textbf{Práctica Experimental de la Unidad IV} de la asignatura
Aplicaciones Web, cuyo resultado de aprendizaje es construir una aplicación web completa bajo el patrón MVC con
un framework moderno, exponiendo e integrando servicios web REST y SOAP, y aplicando criterios de calidad,
seguridad e interoperabilidad. En el marco de la \textbf{evaluación cruzada} solicitada por el docente, el
proyecto evaluado es el \textbf{Sistema de Gestión de Pre-Sustentaciones UTEQ}, desarrollado originalmente por
el Grupo G AMZ y puesto a disposición del Equipo E para su revisión, análisis y retroalimentación técnica.

El sistema tiene como objetivo automatizar la gestión integral de las pre-sustentaciones de trabajos de
titulación en la Universidad Técnica Estatal de Quevedo (UTEQ), cubriendo el ciclo completo: solicitudes de
pre-sustentación, carga de anteproyectos, asignación de jurados, cronogramación de evaluaciones, tutorías,
rúbricas de evaluación, generación de actas y reportes estadísticos. Técnicamente está construido con
\textbf{Spring Boot 3.2.1} (Java 17) en el backend, \textbf{Angular 21.1} en el frontend y
\textbf{PostgreSQL 15} como base de datos, con \textbf{Redis 7} como caché distribuida en Docker Compose. En
las secciones siguientes se desarrolla el fundamento teórico exigido por la guía de práctica (secciones 5.1 a
5.4) y se responden las preguntas de análisis del Anexo C.
